\section{Related Work}
\label{Sec:related}
As mentioned above, \citet{CohenGZ20} demonstrated that function-node splitting triggers rapid convergence of DMS. When applied to standard MS without damping, in most cases it alternates between two solutions of low quality. The DMS-SCFG version converges quickly to a solution of much higher quality. This unexplained success motivated this study.

A general formulation of splitting in the context of belief-propagation was developed by \citet{ruozzi2013}, which showed convergence to local-minima, and some conditions for it to be a global minimum. This was further developed \cite{rebeschini2017} in the specific domain of consensus problems (i.e., as many agents as possible should agree on the same value), in which under several conditions, min-sum using splitting converges to the global optimum in sub-linear time. While formally impressive, these works do not enlighten our understanding of why splitting improves the chances of finding a good outcome in general domains (compared with regular min-sum). Our work, in contrast, wishes to study and understand the \emph{dynamics} of the min-sum mechanism, and how splitting changes these dynamics, ending up with better outcomes.

\citet{DengKL022} identified that the DMS version that uses splitting employs a constant damping parameter for the entire graph. Thus, they proposed an attentive version of the algorithm, in which a deep neural network selects independent damping parameters for every edge of the factor-graph, together with an additional parameter that reduces the weight of incoming messages. With the addition of a specific type of constant splitting, they dramatically improved the quality of solutions produced by DMS-SCFG on benchmarks with random constraints.
This version is not compatible with distributed execution, since the attentive network is trained and applied centrally, with access to the messages of the entire factor graph.% AUTHORS: please verify this characterization of DABP's centralization before submission.
Moreover, as we demonstrate, it loses its advantage on benchmarks with structured constraints. Nevertheless, its success on benchmarks with random constraint costs is striking, and in future work we intend to investigate the patterns of parameter selection that lead to this success and to provide explanations for them.
Most importantly, we demonstrate in our empirical study that the performance of Attentive DMS strongly depends on its use of splitting: without it, the algorithm performs similarly to standard DMS. This is another indication of the importance of investigating the effect of function-node splitting on the performance of DMS, i.e., it further motivates our work.

Our analysis of the two-solution oscillation (Section~\ref{sec:tworoutes}) connects Min-sum on SCFGs to a classical line of results on parallel dynamics with symmetric interactions: synchronous iteration of threshold networks with symmetric weights reaches only fixed points and two-cycles \citep{GolesOlivos80}, as does synchronous weighted-plurality opinion dynamics over finite opinion sets \citep{PoljakSura83}; a period-two result for simultaneous best-response dynamics was recently established for two-player random potential games \citep{AshkenaziGolan25}. These results concern special interaction forms (thresholds, equality indicators, two players); the reduction of committed split Min-sum to synchronous local selection, and the resulting two-solution characterization for arbitrary shared pairwise cost terms, are, to the best of our knowledge, new.
